\chapter{Software Configuration Management}

This chapter documents the Software Configuration Management (SCM) plan for the
AI-Driven Smart Home Energy Optimizer (SHEO) project. SCM encompasses the
principles, techniques, and tools used to identify, control, record, and audit
changes to every work product throughout the project lifecycle. The chapter is
organised around the five canonical SCM tasks drawn from \textbf{IEEE~828}~\cite{ieee828} and
\textbf{IEEE~1042}: configuration identification, version control, change control,
configuration audit, and status reporting.

% ----------------------------------------------------------------
\section{SCM Concepts and Standards}
\label{sec:scm-concepts}
% ----------------------------------------------------------------

\subsection{Foundational Terminology}

Software Configuration Management rests on three foundational concepts that must
be carefully distinguished.

\begin{description}[leftmargin=1.6cm, style=nextline]
  \item[Configuration Item (CI)]
    Any artifact placed under version control and governed by the change-control
    process. A CI is identified by a unique name, a version number, and a
    designated owner. Examples range from source files and test suites to design
    documents and dependency manifests. Once a CI is placed under control, any
    modification to it requires a sanctioned change request.

  \item[Baseline]
    A formally reviewed, approved, and version-locked set of CIs that serves as
    the stable foundation for subsequent work. As defined in \textbf{IEEE~828}, a
    baseline can be changed only through the formal change-control process; ad hoc
    modifications to baselined artifacts are prohibited. Each baseline corresponds
    to a project milestone and is recorded as an immutable, annotated marker in the
    version-control repository.

  \item[Version / Revision / Release]
    These three terms are distinct and must not be conflated. A \textbf{version}
    denotes a particular state of a CI at a point in time, identified by a
    version number of the form \texttt{major.minor.patch} (e.g.\ \texttt{1.0.0}).
    A \textbf{revision} is a specific change to an existing version, typically
    addressing a defect or a minor update within the same major--minor line. A
    \textbf{release} is a version that has been formally approved, baselined, and
    made available for use or delivery; it implies a completed quality gate and is
    associated with a milestone tag.
\end{description}

\subsection{SCM Tasks}

Following \textbf{IEEE~1042}, the SHEO project organises its SCM activity around
five tasks:

\begin{enumerate}[label=\arabic*.]
  \item \textbf{Configuration identification} --- naming and registering every CI
    under control (Section~\ref{sec:scm-cis}).
  \item \textbf{Version control} --- tracking changes to CIs over time using a
    distributed version-control system (Section~\ref{sec:scm-vc}).
  \item \textbf{Change control} --- evaluating, approving, and implementing
    changes to baselined CIs via a defined process
    (Section~\ref{sec:scm-change}).
  \item \textbf{Configuration audit} --- verifying that a baseline is internally
    consistent and satisfies its requirements
    (Section~\ref{sec:scm-audit}).
  \item \textbf{Status reporting} --- recording and communicating the state of
    every CI and change request over time
    (Section~\ref{sec:scm-status}).
\end{enumerate}

% ----------------------------------------------------------------
\section{Configuration Items in SHEO}
\label{sec:scm-cis}
% ----------------------------------------------------------------

The SHEO project identifies seven CI groups. Each group is assigned a short code
used throughout the change-control and status-accounting logs. Table~\ref{tab:cis}
lists the groups with their representative artifacts and designated owners.

\begin{longtable}{@{}
    >{\bfseries}l
    l
    L{4.2cm}
    l
  @{}}
\caption{SHEO Software Configuration Items}
\label{tab:cis} \\
\toprule
\textbf{Code} & \textbf{CI Group} & \textbf{Representative Artifacts} & \textbf{Owner} \\
\midrule
\endfirsthead
\toprule
\textbf{Code} & \textbf{CI Group} & \textbf{Representative Artifacts} & \textbf{Owner} \\
\midrule
\endhead
\bottomrule
\endfoot
C1 & Requirements baseline
   & The approved requirements specification (v1.0) together with its refinement note
   & Requirements lead \\[4pt]
C2 & Design \& engineering docs
   & The software design document, the accompanying structural and behavioural diagrams, and this configuration management plan
   & Architect \\[4pt]
C3 & Backend source
   & The backend service, domain, core, adapter, simulation, persistence, and schema modules together with the data-seeding routine
   & Backend developers \\[4pt]
C4 & Frontend source
   & The client application root, its page and component modules, and its authentication and supporting library code
   & Frontend developers \\[4pt]
C5 & Test suite
   & The shared test fixtures and the nine automated test modules comprising sixty-five cases
   & QA lead \\[4pt]
C6 & Build \& dependency manifests
   & The backend and client dependency and build manifests, with all versions pinned
   & Build owner \\[4pt]
C7 & Runtime configuration
   & The centralised configuration module exposing every requirement-mandated policy threshold, overridable through environment settings
   & Backend lead \\
\bottomrule
\end{longtable}

Several aspects of this identification scheme merit elaboration.

\textbf{Requirements (C1).} The requirements specification is the project's single
source of truth. Functional and non-functional requirements each carry a structured
identifier formed from their category and number. These identifiers serve as the
anchor points for the traceability matrix and are referenced in the recorded
description of every change, maintaining a complete forward and backward trail
between requirements, implementation artifacts, and tests.

\textbf{Runtime configuration (C7).} Every policy threshold mandated by the
requirements is centralised in a single configuration module rather than scattered
as literal constants throughout the source code. The named thresholds --- among them
the offline-detection timeout, the savings-baseline window, the drift tolerance, and
the undo window, each traceable to its governing requirement --- are declared as
named fields whose defaults can be overridden through environment configuration.
Treating policy thresholds as configuration data means that a policy change is a
single, narrowly scoped, reviewable, and baseline-able edit.

\textbf{Build manifests (C6).} The exact version of every backend dependency is
pinned, and the client dependencies are likewise locked to fixed versions. Pinning is
what makes a baseline reproducible: each retrieval of the configuration installs an
identical set of dependency versions.

\textbf{Derived artifacts (non-CIs).} Generated outputs are listed in an ignore
manifest and are not baselined, as they are fully reproducible from the CIs above.
These include the backend virtual environment, compiled-bytecode caches, the
generated runtime database, the installed client dependency tree, and the client
build output.

% ----------------------------------------------------------------
\section{Baselines}
\label{sec:scm-baselines}
% ----------------------------------------------------------------

A baseline freezes a coherent snapshot of all relevant CIs at a milestone boundary.
IEEE~828 requires that once a baseline is established, it may be modified only
through the formal change-control process. In the SHEO project each baseline
corresponds to an immutable, annotated milestone marker placed on the integration
line; the marker records the precise revision it identifies, the date, and the
milestone name.

Table~\ref{tab:baselines} summarises the four project baselines.

\begin{table}[htbp]
\centering
\caption{SHEO Project Baselines}
\label{tab:baselines}
\begin{tabularx}{\textwidth}{@{} l l Y Y @{}}
\toprule
\textbf{Milestone} & \textbf{Marker} & \textbf{CIs Baselined} & \textbf{Entry Criteria} \\
\midrule
M1 --- Requirements
  & M1
  & C1 (the requirements specification v1.0 and its refinement note), the initial documentation scaffold, and the repository with its ignore manifest
  & Requirements specification approved by instructor; all requirement identifiers stable \\[6pt]
M2 --- Design \& skeleton
  & M2
  & C2 (the architectural and structural design documents); the backend module skeleton; C6 with pinned dependencies; C7 with all policy thresholds declared
  & Design reviewed; the skeleton loads cleanly and the backend service starts \\[6pt]
M3 --- Implementation
  & M3
  & C3 (the complete backend services), C4 (all client pages and components), C5 (the full test suite), C7 (all thresholds wired)
  & All sixty-five automated tests pass; the demonstration seeds twenty-one days of data and the application runs end-to-end \\[6pt]
M4 --- Release / submission
  & M4
  & The entire project: C1--C7 consistent end-to-end; the traceability matrix complete; this configuration management plan
  & Test gate passing; traceability matrix fully populated; configuration audit passed \\
\bottomrule
\end{tabularx}
\end{table}

The M1 baseline constitutes the \textbf{Requirements Baseline}: it freezes the
requirements specification and makes subsequent design and implementation decisions
traceable back to an approved set of requirements. The M2 baseline constitutes the
\textbf{Design Baseline}: the software design document and all architectural figures
are locked and serve as the authoritative specification for coding. The M3 and M4 baselines
together constitute the \textbf{Implementation Baseline}: M3 is the first
fully-running system and M4 is the release-ready, fully-audited delivery. Each
baseline approval is recorded in the status-accounting log together with the
milestone marker, the precise revision it identifies, the date, and the list of
change requests incorporated since the previous baseline.

% ----------------------------------------------------------------
\section{Version Control and Branch Workflow}
\label{sec:scm-vc}
% ----------------------------------------------------------------

\subsection{Tooling}

The project uses a single distributed version-control system (Git) to track changes
to every CI over time. The entire project --- source code, tests, documents, and
dependency manifests --- resides in one repository, ensuring that a single baseline
marker captures every CI at once. No CI is managed outside the repository.

\subsection{Branch Model}

A trunk-plus-feature model is adopted, suited to a small, collocated course team
where the integration branch must remain continuously releasable.

\begin{description}[leftmargin=1.6cm, style=nextline]
  \item[Integration line]
    The sole long-lived line of development. It must always be buildable and pass
    the full suite of sixty-five automated tests. Changes are never applied to it
    directly; they arrive exclusively through reviewed merge proposals.
  \item[Feature branch]
    A short-lived line of work that diverges from the integration line for each
    approved change request or new capability. It is discarded once its work has
    been merged.
  \item[Fix branch]
    A short-lived line with the same lifecycle as a feature branch, used
    specifically for defect corrections.
\end{description}

Figure~\ref{fig:branch-tikz} illustrates the branch flow schematically.

\begin{figure}[H]
\centering
\begin{tikzpicture}[
  node distance=1.6cm and 2.2cm,
  commit/.style={circle, draw, fill=gray!20, minimum size=0.55cm, inner sep=0pt,
                 font=\footnotesize},
  tag/.style={rectangle, draw=black!60, rounded corners=2pt, fill=yellow!20,
              font=\tiny, inner sep=2pt},
  branch/.style={font=\small\itshape},
  arr/.style={->, >=stealth, thick}
]
%% main branch (horizontal)
\node[commit] (c0) {};
\node[commit, right=of c0] (c1) {};
\node[commit, right=of c1] (c2) {};
\node[commit, right=of c2] (c3) {};
\draw[arr] (c0) -- (c1);
\draw[arr] (c1) -- (c2);
\draw[arr] (c2) -- (c3);
%% feature branch (above)
\node[commit, above right=0.9cm and 0.6cm of c0] (f1) {};
\node[commit, right=1.4cm of f1] (f2) {};
\draw[arr] (c0) -- (f1);
\draw[arr] (f1) -- (f2);
\draw[arr] (f2) -- (c2);
%% fix branch (below)
\node[commit, below right=0.9cm and 0.6cm of c1] (b1) {};
\draw[arr] (c1) -- (b1);
\draw[arr] (b1) -- (c3);
%% tags
\node[tag, above=0.22cm of c0] {M1};
\node[tag, above=0.22cm of c3] {M2};
%% labels
\node[branch, left=0.3cm of c0] {integration};
\node[branch, above=0.1cm of f1] {feature};
\node[branch, below=0.1cm of b1] {fix};
\end{tikzpicture}
\caption{Schematic branch model: feature and fix branches diverge from the
         integration line and are merged back after peer review and a passing test
         gate. Milestone markers (M1, M2, \ldots) denote baseline revisions on the
         integration line.}
\label{fig:branch-tikz}
\end{figure}

\subsection{Commit Conventions}

The recorded description of each change follows a structured convention so that the
project history remains scannable and every change stays traceable to its origin.
Each description pairs a change-type --- distinguishing, for instance, a new
capability, a defect correction, an addition to the test suite, a documentation
update, an internal restructuring, or routine maintenance --- with the affected
subsystem and, where applicable, the identifier of the requirement it implements or
modifies.\footnote{For example, a change introducing conflict detection among
enabled rules would be recorded as a new capability scoped to the rules subsystem and
annotated with the functional requirement it satisfies.} Naming the governing
requirement identifier is mandatory whenever a change implements or modifies a
requirement, as this is what keeps automated status accounting tractable.

\subsection{Version Numbering}

Software artifacts are versioned according to a three-part
\emph{major.minor.patch} scheme, with version \texttt{1.0.0} designating the M4
release. The \textbf{major} digit increments on incompatible interface or data-model
changes; the \textbf{minor} digit increments on backward-compatible new
functionality; the \textbf{patch} digit increments on backward-compatible defect
fixes. The backend declares its own version number, which is advanced as part of the
same change that establishes a new baseline.

% ----------------------------------------------------------------
\section{Change Control and Status Reporting}
\label{sec:scm-change}
% ----------------------------------------------------------------

\subsection{Change Control Process}

Any modification to a baselined CI --- whether it addresses a new requirement, a
discovered defect, an adjustment to a policy threshold, or a document revision ---
must pass through a three-stage change-control process. For a course-sized team this
process is intentionally lightweight, but the evaluation and review gates are real
and enforced.

\textbf{Stage~1: Request, Evaluate, and Approve.}
A team member raises a tracked request describing the proposed change, the CIs
affected (by code C1--C7), and the requirement identifiers involved. A designated
reviewer assesses the impact: which requirements are affected, which modules must
change, which tests must be updated, and --- critically --- whether a baselined CI
is implicated (changes to a baseline require stronger justification than changes to
in-progress work). The configuration-control authority approves, defers, or rejects
the request and records the decision against it.

\textbf{Stage~2: Implement on Branch and Peer-Review.}
An approved request is implemented on a dedicated short-lived feature or fix branch
that diverges from the integration line. The implementing developer records the code
change together with the tests that demonstrate its correctness, naming the relevant
requirement identifier in each change description. When implementation is complete, a
merge proposal is opened and a second team member performs a formal peer review,
verifying correctness, adherence to the three-tier layered architecture, and that
every change is traceable to an approved requirement. Self-merges are not permitted.

\textbf{Stage~3: Test Gate, Merge, and Release.}
The full automated test suite of sixty-five cases must pass on the branch before
merge. This is the non-negotiable merge gate. On passing, the branch is merged into
the integration line and then discarded. If the merged change completes a milestone,
the configuration manager establishes the corresponding annotated baseline marker
(M1--M4) and updates the status-accounting log and traceability matrix.

\subsection{Configuration Status Reporting}
\label{sec:scm-status}

Configuration status accounting answers three questions for every CI at any point
in time: \emph{what} changed, \emph{who} made the change, and \emph{when} it was
made. In the SHEO project, status accounting is maintained through two artefacts.

The \textbf{change register} records every change request and defect with a unique
identifier, the affected CI codes, the requirement identifiers involved, the current
state as it advances from raised, through in progress and under review, to merged and
finally released, and the line of work and merge proposal that resolved it. The
change identifier is carried through to the line of work and the recorded change
descriptions, providing a complete forward and backward audit trail.

The \textbf{baseline log} records, for each milestone, its marker, the precise
revision it identifies, the calendar date, and the list of change identifiers
incorporated since the preceding baseline. This log is the canonical record of what
the project delivered at each stage and forms the primary input to the configuration
audit.

\subsection{Configuration Audit}
\label{sec:scm-audit}

Two complementary audits are performed at each milestone, and most rigorously at
the M4 release baseline:

\begin{description}[leftmargin=1.6cm, style=nextline]
  \item[Functional Configuration Audit (FCA)]
    Verifies that the baseline satisfies all requirements assigned to it. Evidence
    is the requirements-to-implementation traceability matrix, which maps each
    requirement identifier to the implementing module and routine and to the test
    case that exercises it, combined with a fully passing run of the test suite.

  \item[Physical Configuration Audit (PCA)]
    Verifies that the baseline contains exactly the expected CIs (C1--C7) at their
    declared versions, and that all derived artifacts are absent (governed by the
    ignore manifest). The auditor checks the contents identified by the baseline
    marker against the CI catalogue in Section~\ref{sec:scm-cis}.
\end{description}

Together these two audits provide the assurance required by IEEE~828 that the
released configuration is both correct and complete, closing the SCM lifecycle for
each milestone.
